\documentclass{article}
\usepackage{amsmath, amssymb, amsthm}
\usepackage{hyperref}
\usepackage{geometry}
\geometry{margin=1in}

\title{Erd\H{o}s Problem \#1133}
\date{Last updated: 2026-01-01}
\author{Source: erdosproblems.com}

\begin{document}
\maketitle

\noindent\textbf{Status:} OPEN \quad \textbf{No prize}

\noindent\textbf{Tags:} analysis, polynomials

\medskip
\noindent\textbf{Problem Statement:}

Let $C>0$. There exists $\epsilon>0$ such that if $n$ is sufficiently large the following holds.

For any $x_1,\ldots,x_n\in [-1,1]$ there exist $y_1,\ldots,y_n\in [-1,1]$ such that, if $P$ is a polynomial of degree $m<(1+\epsilon)n$ with $P(x_i)=y_i$ for at least $(1-\epsilon)n$ many $1\leq i\leq n$, then\[\max_{x\in [-1,1]}\lvert P(x)\rvert >C.\]

\bigskip
\noindent\textbf{Remarks:}

Erd\H{o}s proved that, for any $C>0$, there exists $\epsilon>0$ such that if $n$ is sufficiently large and $m=\lfloor (1+\epsilon)n\rfloor$ then for any $x_1,\ldots,x_m\in [-1,1]$ there is a polynomial $P$ of degree $n$ such that $\lvert P(x_i)\rvert\leq 1$ for $1\leq i\leq m$ and\[\max_{x\in [-1,1]}\lvert P(x)\rvert>C.\]The conjectured statement would also imply this, but Erd\H{o}s in \cite{Er67} says he could not even prove it for $m=n$.

\bigskip
\noindent\textbf{References:}

[Er67] Erd\H{o}s, P., Problems and results on the convergence and divergence properties of the Lagrange interpolation polynomials and some extremal problems. Mathematica (Cluj) (1967), 65-73.

\bigskip
\noindent\small{Source: \url{https://www.erdosproblems.com/1133}}
\end{document}
